% ============================================================================
% Configuration LaTeX pour des codes bien présentés avec couleurs (style IDE)
% À inclure dans le préambule du document principal avec : % ============================================================================
% Configuration LaTeX pour des codes bien présentés avec couleurs (style IDE)
% À inclure dans le préambule du document principal avec : % ============================================================================
% Configuration LaTeX pour des codes bien présentés avec couleurs (style IDE)
% À inclure dans le préambule du document principal avec : % ============================================================================
% Configuration LaTeX pour des codes bien présentés avec couleurs (style IDE)
% À inclure dans le préambule du document principal avec : \input{code_styles}
% ============================================================================

\usepackage{listings}
\usepackage{xcolor}

% Définition des couleurs (style VS Code / PyCharm)
\definecolor{codegreen}{rgb}{0,0.6,0}
\definecolor{codegray}{rgb}{0.5,0.5,0.5}
\definecolor{codepurple}{rgb}{0.58,0,0.82}
\definecolor{codeblue}{rgb}{0.13,0.29,0.53}
\definecolor{codeorange}{rgb}{0.8,0.4,0}
\definecolor{backcolour}{rgb}{0.95,0.95,0.92}
\definecolor{framecolor}{rgb}{0.7,0.7,0.7}

% Style général pour Python
\lstdefinestyle{pythoncode}{
    backgroundcolor=\color{backcolour},
    commentstyle=\color{codegreen}\itshape,
    keywordstyle=\color{codeblue}\bfseries,
    numberstyle=\tiny\color{codegray},
    stringstyle=\color{codeorange},
    basicstyle=\ttfamily\small,
    breakatwhitespace=false,
    breaklines=true,
    captionpos=b,
    keepspaces=true,
    numbers=left,
    numbersep=8pt,
    showspaces=false,
    showstringspaces=false,
    showtabs=false,
    tabsize=4,
    frame=single,
    framerule=0.8pt,
    rulecolor=\color{framecolor},
    xleftmargin=15pt,
    xrightmargin=5pt,
    framexleftmargin=12pt,
    aboveskip=15pt,
    belowskip=10pt
}

% Style pour YAML (Docker Compose)
\lstdefinestyle{yamlcode}{
    backgroundcolor=\color{backcolour},
    commentstyle=\color{codegreen}\itshape,
    keywordstyle=\color{codeblue}\bfseries,
    numberstyle=\tiny\color{codegray},
    stringstyle=\color{codeorange},
    basicstyle=\ttfamily\small,
    breakatwhitespace=false,
    breaklines=true,
    captionpos=b,
    keepspaces=true,
    numbers=left,
    numbersep=8pt,
    showspaces=false,
    showstringspaces=false,
    showtabs=false,
    tabsize=2,
    frame=single,
    framerule=0.8pt,
    rulecolor=\color{framecolor},
    xleftmargin=15pt,
    xrightmargin=5pt,
    framexleftmargin=12pt,
    aboveskip=15pt,
    belowskip=10pt
}

% Style pour Bash/Shell
\lstdefinestyle{bashcode}{
    backgroundcolor=\color{backcolour},
    commentstyle=\color{codegreen}\itshape,
    keywordstyle=\color{codeblue}\bfseries,
    numberstyle=\tiny\color{codegray},
    stringstyle=\color{codeorange},
    basicstyle=\ttfamily\small,
    breakatwhitespace=false,
    breaklines=true,
    captionpos=b,
    keepspaces=true,
    numbers=left,
    numbersep=8pt,
    showspaces=false,
    showstringspaces=false,
    showtabs=false,
    tabsize=4,
    frame=single,
    framerule=0.8pt,
    rulecolor=\color{framecolor},
    xleftmargin=15pt,
    xrightmargin=5pt,
    framexleftmargin=12pt,
    aboveskip=15pt,
    belowskip=10pt
}

% Configuration par défaut selon le langage
\lstset{
    language=Python,
    style=pythoncode
}

% Mots-clés supplémentaires pour Python
\lstset{
    morekeywords={self, True, False, None, as, with, async, await},
    emph={torch, nn, optim, F, np, pd, plt, sklearn},
    emphstyle=\color{codepurple}\bfseries
}

% Instructions d'utilisation :
% Pour Python : \begin{lstlisting}[language=Python, caption=Titre du code]
% Pour YAML   : \begin{lstlisting}[language=yaml, style=yamlcode, caption=Titre]
% Pour Bash   : \begin{lstlisting}[language=bash, style=bashcode, caption=Titre]

% ============================================================================

\usepackage{listings}
\usepackage{xcolor}

% Définition des couleurs (style VS Code / PyCharm)
\definecolor{codegreen}{rgb}{0,0.6,0}
\definecolor{codegray}{rgb}{0.5,0.5,0.5}
\definecolor{codepurple}{rgb}{0.58,0,0.82}
\definecolor{codeblue}{rgb}{0.13,0.29,0.53}
\definecolor{codeorange}{rgb}{0.8,0.4,0}
\definecolor{backcolour}{rgb}{0.95,0.95,0.92}
\definecolor{framecolor}{rgb}{0.7,0.7,0.7}

% Style général pour Python
\lstdefinestyle{pythoncode}{
    backgroundcolor=\color{backcolour},
    commentstyle=\color{codegreen}\itshape,
    keywordstyle=\color{codeblue}\bfseries,
    numberstyle=\tiny\color{codegray},
    stringstyle=\color{codeorange},
    basicstyle=\ttfamily\small,
    breakatwhitespace=false,
    breaklines=true,
    captionpos=b,
    keepspaces=true,
    numbers=left,
    numbersep=8pt,
    showspaces=false,
    showstringspaces=false,
    showtabs=false,
    tabsize=4,
    frame=single,
    framerule=0.8pt,
    rulecolor=\color{framecolor},
    xleftmargin=15pt,
    xrightmargin=5pt,
    framexleftmargin=12pt,
    aboveskip=15pt,
    belowskip=10pt
}

% Style pour YAML (Docker Compose)
\lstdefinestyle{yamlcode}{
    backgroundcolor=\color{backcolour},
    commentstyle=\color{codegreen}\itshape,
    keywordstyle=\color{codeblue}\bfseries,
    numberstyle=\tiny\color{codegray},
    stringstyle=\color{codeorange},
    basicstyle=\ttfamily\small,
    breakatwhitespace=false,
    breaklines=true,
    captionpos=b,
    keepspaces=true,
    numbers=left,
    numbersep=8pt,
    showspaces=false,
    showstringspaces=false,
    showtabs=false,
    tabsize=2,
    frame=single,
    framerule=0.8pt,
    rulecolor=\color{framecolor},
    xleftmargin=15pt,
    xrightmargin=5pt,
    framexleftmargin=12pt,
    aboveskip=15pt,
    belowskip=10pt
}

% Style pour Bash/Shell
\lstdefinestyle{bashcode}{
    backgroundcolor=\color{backcolour},
    commentstyle=\color{codegreen}\itshape,
    keywordstyle=\color{codeblue}\bfseries,
    numberstyle=\tiny\color{codegray},
    stringstyle=\color{codeorange},
    basicstyle=\ttfamily\small,
    breakatwhitespace=false,
    breaklines=true,
    captionpos=b,
    keepspaces=true,
    numbers=left,
    numbersep=8pt,
    showspaces=false,
    showstringspaces=false,
    showtabs=false,
    tabsize=4,
    frame=single,
    framerule=0.8pt,
    rulecolor=\color{framecolor},
    xleftmargin=15pt,
    xrightmargin=5pt,
    framexleftmargin=12pt,
    aboveskip=15pt,
    belowskip=10pt
}

% Configuration par défaut selon le langage
\lstset{
    language=Python,
    style=pythoncode
}

% Mots-clés supplémentaires pour Python
\lstset{
    morekeywords={self, True, False, None, as, with, async, await},
    emph={torch, nn, optim, F, np, pd, plt, sklearn},
    emphstyle=\color{codepurple}\bfseries
}

% Instructions d'utilisation :
% Pour Python : \begin{lstlisting}[language=Python, caption=Titre du code]
% Pour YAML   : \begin{lstlisting}[language=yaml, style=yamlcode, caption=Titre]
% Pour Bash   : \begin{lstlisting}[language=bash, style=bashcode, caption=Titre]

% ============================================================================

\usepackage{listings}
\usepackage{xcolor}

% Définition des couleurs (style VS Code / PyCharm)
\definecolor{codegreen}{rgb}{0,0.6,0}
\definecolor{codegray}{rgb}{0.5,0.5,0.5}
\definecolor{codepurple}{rgb}{0.58,0,0.82}
\definecolor{codeblue}{rgb}{0.13,0.29,0.53}
\definecolor{codeorange}{rgb}{0.8,0.4,0}
\definecolor{backcolour}{rgb}{0.95,0.95,0.92}
\definecolor{framecolor}{rgb}{0.7,0.7,0.7}

% Style général pour Python
\lstdefinestyle{pythoncode}{
    backgroundcolor=\color{backcolour},
    commentstyle=\color{codegreen}\itshape,
    keywordstyle=\color{codeblue}\bfseries,
    numberstyle=\tiny\color{codegray},
    stringstyle=\color{codeorange},
    basicstyle=\ttfamily\small,
    breakatwhitespace=false,
    breaklines=true,
    captionpos=b,
    keepspaces=true,
    numbers=left,
    numbersep=8pt,
    showspaces=false,
    showstringspaces=false,
    showtabs=false,
    tabsize=4,
    frame=single,
    framerule=0.8pt,
    rulecolor=\color{framecolor},
    xleftmargin=15pt,
    xrightmargin=5pt,
    framexleftmargin=12pt,
    aboveskip=15pt,
    belowskip=10pt
}

% Style pour YAML (Docker Compose)
\lstdefinestyle{yamlcode}{
    backgroundcolor=\color{backcolour},
    commentstyle=\color{codegreen}\itshape,
    keywordstyle=\color{codeblue}\bfseries,
    numberstyle=\tiny\color{codegray},
    stringstyle=\color{codeorange},
    basicstyle=\ttfamily\small,
    breakatwhitespace=false,
    breaklines=true,
    captionpos=b,
    keepspaces=true,
    numbers=left,
    numbersep=8pt,
    showspaces=false,
    showstringspaces=false,
    showtabs=false,
    tabsize=2,
    frame=single,
    framerule=0.8pt,
    rulecolor=\color{framecolor},
    xleftmargin=15pt,
    xrightmargin=5pt,
    framexleftmargin=12pt,
    aboveskip=15pt,
    belowskip=10pt
}

% Style pour Bash/Shell
\lstdefinestyle{bashcode}{
    backgroundcolor=\color{backcolour},
    commentstyle=\color{codegreen}\itshape,
    keywordstyle=\color{codeblue}\bfseries,
    numberstyle=\tiny\color{codegray},
    stringstyle=\color{codeorange},
    basicstyle=\ttfamily\small,
    breakatwhitespace=false,
    breaklines=true,
    captionpos=b,
    keepspaces=true,
    numbers=left,
    numbersep=8pt,
    showspaces=false,
    showstringspaces=false,
    showtabs=false,
    tabsize=4,
    frame=single,
    framerule=0.8pt,
    rulecolor=\color{framecolor},
    xleftmargin=15pt,
    xrightmargin=5pt,
    framexleftmargin=12pt,
    aboveskip=15pt,
    belowskip=10pt
}

% Configuration par défaut selon le langage
\lstset{
    language=Python,
    style=pythoncode
}

% Mots-clés supplémentaires pour Python
\lstset{
    morekeywords={self, True, False, None, as, with, async, await},
    emph={torch, nn, optim, F, np, pd, plt, sklearn},
    emphstyle=\color{codepurple}\bfseries
}

% Instructions d'utilisation :
% Pour Python : \begin{lstlisting}[language=Python, caption=Titre du code]
% Pour YAML   : \begin{lstlisting}[language=yaml, style=yamlcode, caption=Titre]
% Pour Bash   : \begin{lstlisting}[language=bash, style=bashcode, caption=Titre]

% ============================================================================

\usepackage{listings}
\usepackage{xcolor}

% Définition des couleurs (style VS Code / PyCharm)
\definecolor{codegreen}{rgb}{0,0.6,0}
\definecolor{codegray}{rgb}{0.5,0.5,0.5}
\definecolor{codepurple}{rgb}{0.58,0,0.82}
\definecolor{codeblue}{rgb}{0.13,0.29,0.53}
\definecolor{codeorange}{rgb}{0.8,0.4,0}
\definecolor{backcolour}{rgb}{0.95,0.95,0.92}
\definecolor{framecolor}{rgb}{0.7,0.7,0.7}

% Style général pour Python
\lstdefinestyle{pythoncode}{
    backgroundcolor=\color{backcolour},
    commentstyle=\color{codegreen}\itshape,
    keywordstyle=\color{codeblue}\bfseries,
    numberstyle=\tiny\color{codegray},
    stringstyle=\color{codeorange},
    basicstyle=\ttfamily\small,
    breakatwhitespace=false,
    breaklines=true,
    captionpos=b,
    keepspaces=true,
    numbers=left,
    numbersep=8pt,
    showspaces=false,
    showstringspaces=false,
    showtabs=false,
    tabsize=4,
    frame=single,
    framerule=0.8pt,
    rulecolor=\color{framecolor},
    xleftmargin=15pt,
    xrightmargin=5pt,
    framexleftmargin=12pt,
    aboveskip=15pt,
    belowskip=10pt
}

% Style pour YAML (Docker Compose)
\lstdefinestyle{yamlcode}{
    backgroundcolor=\color{backcolour},
    commentstyle=\color{codegreen}\itshape,
    keywordstyle=\color{codeblue}\bfseries,
    numberstyle=\tiny\color{codegray},
    stringstyle=\color{codeorange},
    basicstyle=\ttfamily\small,
    breakatwhitespace=false,
    breaklines=true,
    captionpos=b,
    keepspaces=true,
    numbers=left,
    numbersep=8pt,
    showspaces=false,
    showstringspaces=false,
    showtabs=false,
    tabsize=2,
    frame=single,
    framerule=0.8pt,
    rulecolor=\color{framecolor},
    xleftmargin=15pt,
    xrightmargin=5pt,
    framexleftmargin=12pt,
    aboveskip=15pt,
    belowskip=10pt
}

% Style pour Bash/Shell
\lstdefinestyle{bashcode}{
    backgroundcolor=\color{backcolour},
    commentstyle=\color{codegreen}\itshape,
    keywordstyle=\color{codeblue}\bfseries,
    numberstyle=\tiny\color{codegray},
    stringstyle=\color{codeorange},
    basicstyle=\ttfamily\small,
    breakatwhitespace=false,
    breaklines=true,
    captionpos=b,
    keepspaces=true,
    numbers=left,
    numbersep=8pt,
    showspaces=false,
    showstringspaces=false,
    showtabs=false,
    tabsize=4,
    frame=single,
    framerule=0.8pt,
    rulecolor=\color{framecolor},
    xleftmargin=15pt,
    xrightmargin=5pt,
    framexleftmargin=12pt,
    aboveskip=15pt,
    belowskip=10pt
}

% Configuration par défaut selon le langage
\lstset{
    language=Python,
    style=pythoncode
}

% Mots-clés supplémentaires pour Python
\lstset{
    morekeywords={self, True, False, None, as, with, async, await},
    emph={torch, nn, optim, F, np, pd, plt, sklearn},
    emphstyle=\color{codepurple}\bfseries
}

% Instructions d'utilisation :
% Pour Python : \begin{lstlisting}[language=Python, caption=Titre du code]
% Pour YAML   : \begin{lstlisting}[language=yaml, style=yamlcode, caption=Titre]
% Pour Bash   : \begin{lstlisting}[language=bash, style=bashcode, caption=Titre]
